% 第10章 自由伝播と視線積分
\section{自由伝播と視線積分}
\label{ch:ch10_los}

\paragraph{この章で導くこと（入力→出力）}
ボルツマン方程式（入力）の形式解＝視線（line-of-sight）積分を導き、ソース関数を
SW・ISW・Doppler・四重極補正の4項に分解する（出力 S10.1--S10.3）。これにより
「最終散乱面のスナップショット＋道中の積分効果」という描像が数式になり、
第\ref{ch:ch11_cl}章の $C_\ell$ 公式の被積分核 $\Theta_\ell(k)$ が定まる。

\subsection{形式解（導出契約 S10.1）}
第\ref{ch:ch06_boltzmann}章の光子方程式は、$\mu=\hat k\cdot\hat p$ として
\begin{equation}
  \Theta'+\frac{ck\mu}{\mathcal H}\Theta
  =-\Phi'-\frac{ck\mu}{\mathcal H}\Psi-\tau'\big[\Theta_0-\Theta+\mu v_b-\tfrac12 P_2(\mu)\Theta_2\big],
\end{equation}
と書ける。左辺を積分因子 $e^{ik\mu(\eta-\eta_0)-\tau}$ でまとめると厳密解
\begin{equation}
  \Theta(\eta_0,k,\mu)=\int_{-\infty}^{0}\tilde S(k,x)\,e^{ik\mu(\eta(x)-\eta_0)}\,dx,
\end{equation}
を得る（$\tilde S$ は下記）。観測されるのは $\mu$ 方向の多重極だから、次に $\mu$ 依存を
ベッセル関数に変換する。

\subsection{$\mu$ の部分積分と $j_\ell$ の出現（導出契約 S10.2）}
平面波の Rayleigh 展開 $e^{ik\mu\Delta\eta}=\sum_\ell(2\ell+1)i^\ell P_\ell(\mu)
j_\ell(k\Delta\eta)$ を使い、$\mu$ の冪を $\frac{1}{ik}\frac{d}{d\Delta\eta}$ への
置換（部分積分）で吸収すると、温度多重極は
\begin{equation}
  \boxed{\ \Theta_\ell(\eta_0,k)=\int_{-\infty}^{0}\tilde S(k,x)\,
  j_\ell\big[k(\eta_0-\eta(x))\big]\,dx\ }
  \label{eq:los}
\end{equation}
となる。$j_\ell[k(\eta_0-\eta)]$ は「波数 $k$ の擾乱を角度スケール $\ell$ へ射影する核」
であり、$\ell\approx k(\eta_0-\eta_*)=k\chi_*$ のとき最大の寄与を与える。

\subsection{ソース関数の物理}
4項に分解されたソース関数は（偏光なし $\Pi=\Theta_2$）
\begin{align}
  \tilde S(k,x)=&\ \underbrace{\tilde g\Big[\Theta_0+\Psi+\tfrac14\Theta_2\Big]}_{\text{SW（実効温度）}}
  +\underbrace{e^{-\tau}\big[\Psi'-\Phi'\big]}_{\text{ISW}} \notag\\
  &-\underbrace{\frac{1}{ck}\frac{d}{dx}\big(\mathcal H\tilde g\,v_b\big)}_{\text{Doppler}}
  +\underbrace{\frac{3}{4c^2k^2}\frac{d}{dx}\!\Big[\mathcal H\frac{d}{dx}\big(\mathcal H\tilde g\,\Theta_2\big)\Big]}_{\text{四重極補正}}.
  \label{eq:source}
\end{align}
\begin{itemize}
  \item \textbf{SW項}: $\tilde g$ が最終散乱面に重みを与え、そこでの\emph{実効温度}
  $\Theta_0+\Psi$ を射影する。$\Psi$ はポテンシャル井戸からの脱出赤方偏移。
  \item \textbf{ISW項}: $e^{-\tau}$（再結合後はほぼ1）と $\Psi'-\Phi'$ の積。道中で
  ポテンシャルが時間変化すると光子が正味のエネルギーを得る（第\ref{ch:ch12_peaks}章）。
  \item \textbf{Doppler項}: バリオン速度 $v_b$ の視線方向成分。
  \item \textbf{四重極補正}: 偏光を無視した近似での残差項。
\end{itemize}

\begin{numreality}
式\eqref{eq:source} は \texttt{cmbcore.spectrum.\_source} に4項個別に実装され、
微分はスプライン経由で評価される。式\eqref{eq:los} の数値積分には注意が必要で、
$\Theta_\ell(k)$ は $k$ について周期 $\sim2\pi/\chi_*$ で急速振動する。そこで
\textbf{ソース $\tilde S(k,x)$ を粗い $k$ で解いて $k$ 方向にスプライン補間し、
細かい $k$ グリッドで式\eqref{eq:los} を実行}する（CLASS/CAMB と同じ手法、
第\ref{ch:ch14_numerics}章）。
\end{numreality}

\subsection{瞬間再結合極限と Sachs-Wolfe（導出契約 S10.3）}
$\tilde g\to\delta(x-x_*)$ の極限で SW 項は
$\Theta_\ell\simeq[\Theta_0+\Psi](x_*)\,j_\ell(k\chi_*)$ となる。超地平線・物質優勢では
第\ref{ch:ch07_initial}章の初期条件と組み合わせて
\begin{equation}
  \Theta_0+\Psi=\tfrac13\Psi,
\end{equation}
すなわち「温度ゆらぎはポテンシャルの $1/3$」という有名な Sachs-Wolfe 公式を得る。
これが低 $\ell$ プラトーの解析的基礎（第\ref{ch:ch12_peaks}章）になる。

\paragraph{まとめ・チェックリスト}
\begin{itemize}
  \item 形式解→部分積分で式\eqref{eq:los}（$\Theta_\ell=\int\tilde S\,j_\ell\,dx$）を導いた。
  \item ソースを SW/ISW/Doppler/四重極の4項に分解した\eqref{eq:source}。
  \item 数値核 $\Theta_\ell(k)$ の急速 $k$ 振動と、その解像法を明示した。
\end{itemize}

\paragraph{演習}
(1) Rayleigh 展開から式\eqref{eq:los} を $\ell=2$ について書き下せ。
(2) NB4 でソース4項を分解し、各項が効く $x$ 領域を読み取れ。
(3) SW 極限 $\Theta_0+\Psi=\Psi/3$ を初期条件から確認せよ。
